\documentclass[11pt]{article}
\usepackage[T1]{fontenc}
\usepackage[margin=1in]{geometry}
\usepackage{amsmath,amssymb,amsthm,mathtools,booktabs}
\usepackage[colorlinks=true,linkcolor=blue,urlcolor=blue,citecolor=blue]{hyperref}
\usepackage{microtype}
\newtheorem{theorem}{Theorem}[section]
\newtheorem{lemma}[theorem]{Lemma}
\newtheorem{proposition}[theorem]{Proposition}
\newtheorem{corollary}[theorem]{Corollary}
\theoremstyle{definition}
\newtheorem{definition}[theorem]{Definition}
\theoremstyle{remark}
\newtheorem{remark}[theorem]{Remark}
\newcommand{\one}{\mathbf1}
\newcommand{\E}{\mathbb E}
\newcommand{\dd}{\,\mathrm d}
\newcommand{\Tr}{\operatorname{Tr}}
\newcommand{\SU}{\mathrm{SU}(2)}
\newcommand{\norm}[1]{\left\lVert#1\right\rVert}
\newcommand{\supp}{\operatorname{supp}}
\setlength{\emergencystretch}{3em}
\hypersetup{pdftitle={Exact sparse Wilson marginals and one step repair bounds},pdfauthor={Research note},pdfsubject={Finite SU(2) lattice marginalization, not an RG flow or continuum theorem}}
\title{Exact sparse Wilson marginals and one step repair bounds}
\author{Research note}
\date{October 9, 2026}
\begin{document}
\maketitle
\begin{abstract}
We integrate one $SU(2)$ Wilson link, then any fixed family of links no
two of which share a plaquette, exactly. Each integrated star contributes
$e^{-6H}F(Hb)$, where $b$ is the norm of its six-staple sum and
$F(t)=2I_1(t)/t$, with its continuous value one at zero. Extracting the
conditional minimum $6-b$ leaves a bounded remainder of size $O(\log H)$.
The resulting interaction is genuinely not a sum of original
elementary-cell interactions, so the preceding weighted class is not
closed even under one removal. Nevertheless, a retained-link edit lifts
to at most seven times as many original links. The original Wilson
energy moment and minimizer-cap estimate give a one-region weighted
repair/exception bound in the exact marginal, with explicit local costs.
For a reflection-invariant sparse pattern the complete marginal inherits
site-reflection positivity; no positivity of its individual induced terms
is assumed. The rough labels after integration are declared separately
from missing original plaquettes. This is one exact partial
marginalization, with no rescaling or repeated-RG stability, no signed
covering construction, and no continuum or mass-gap conclusion. No
novelty claim is made.
\end{abstract}
\tableofcontents

\section{Model and the removed-link set}
We start with the \emph{pure} Wilson model, not with an arbitrary weighted
endpoint from the fourth pass. On
$\Lambda_L=(\mathbb Z/L\mathbb Z)^4$, $L=2^m$, $m\ge2$, set
\[
 E(U)=\sum_qs(U_q),\quad s(g)=1-\tfrac12\Tr g,
 \quad\mu_H(\dd U)=Z_H^{-1}e^{-HE(U)}\dd U,\quad H\ge1.
\]
Every positively oriented link has a variable in $\SU$; its reverse has
the inverse. Each unoriented plaquette is counted once. The reference
$\dd U$ is product \emph{probability} Haar. Let $D$ be a fixed set of
removed links such that each plaquette contains at most one link of $D$.
Write $R$ for the retained links and $z=U_R$. The exact marginal is
$\nu_D=(U\mapsto U_R)_*\mu_H$. Selection of $D$ from the source
configuration is not part of this statement.

For $e\in D$, its six plaquettes have energy
\[
 E_e(u_e,z)=\sum_{j=1}^6\left(1-\tfrac12\Tr(u_eS_{ej}(z)^{-1})\right),
\]
where $S_{ej}$ is the three-link detour from the source of $e$ to its
target. All detour links are retained. Identify these matrices with unit
quaternions, put
\begin{equation}
 B_e=\sum_{j=1}^6s_{ej},\qquad b_e=|B_e|\in[0,6],\qquad m_e=6-b_e.
 \label{eq:staples}
\end{equation}
The quaternion inner product gives $E_e=6-u_e\cdot B_e$ and
$m_e=\min_{u_e}E_e$. For $b_e=0$ every $u_e$ is a minimizer. Gauge
transformations act by common left/right endpoint multiplication on the
staples, preserving $b_e$ and $m_e$.

Let $E_0(z)$ sum exactly those original plaquettes containing no removed
link. Since the removed stars have disjoint plaquette sets,
\begin{equation}
 E(U_D,z)=E_0(z)+\sum_{e\in D}E_e(u_e,z).
 \label{eq:decomposition}
\end{equation}
Stars may share retained links and vertices. Conditional factorization
in the removed variables does not mean independence of the induced
factors under the marginal law.

\section{Exact integral and a uniform logarithmic remainder}
Define for $t\ge0$
\begin{equation}
 F(t)=\frac2\pi\int_0^\pi e^{t\cos\theta}\sin^2\theta\dd\theta
     =\sum_{j=0}^\infty\frac{(t^2/4)^j}{j!(j+1)!}.
 \label{eq:F}
\end{equation}
Thus $F(0)=1$ and $F(t)=2I_1(t)/t$ for $t>0$. This normalization agrees
with the primary integral and power-series definitions
\cite[10.32.2, 10.25.2]{dlmf}. The series can also be obtained directly
by expanding the exponential and integrating the even spherical moments.

\begin{theorem}[Exact retained density]\label{thm:marginal}
For every admissible fixed $D$,
\begin{equation}
 \nu_D(\dd z)=\frac1{Z_H}e^{-HE_0(z)}
        \prod_{e\in D}\left[e^{-6H}F(Hb_e(z))\right]\dd z.
 \label{eq:exact}
\end{equation}
There is no extra gauge-volume factor or unknown normalizing constant:
integrating the displayed numerator gives exactly $Z_H$.
\end{theorem}
\begin{proof}
For $b_e>0$, left multiplication by the unit quaternion $B_e/b_e$
rotates the linear form $u_e\cdot B_e$ to $b_e\cos\theta$. Probability
Haar on $\SU\simeq S^3$ has angular density
$(2/\pi)\sin^2\theta\dd\theta$. The integral over $u_e$ is therefore
$e^{-6H}F(Hb_e)$. At $b_e=0$ it is $e^{-6H}$, in agreement with
$F(0)=1$. Apply Fubini to \eqref{eq:decomposition}. All factors are
positive, bounded and measurable in this finite volume.
\end{proof}

Set
\begin{equation}
 r(t)=t-\log F(t),\qquad
 r_H=\tfrac12+\log24+\tfrac32\log(1+6H).
 \label{eq:r}
\end{equation}
\begin{lemma}[Uniform remainder including zero staple sum]\label{lem:r}
For every $t\ge0$,
\begin{equation}
 0\le r(t)\le\tfrac12+\log24+\tfrac32\log(1+t),\quad r(0)=0.
 \label{eq:rbound}
\end{equation}
In particular $0\le r(Hb_e)\le r_H$. For $t>0$,
$0<F'(t)/F(t)<1$, so $0<r'(t)<1$.
\end{lemma}
\begin{proof}
Write
$e^{-r(t)}=\E_{\mathrm{Haar}}e^{-t(1-\cos\theta)}\le1$.
On the cap $\theta\le\rho=(1+t)^{-1/2}\le1$,
$t(1-\cos\theta)\le t\rho^2/2\le1/2$.
The previously proved elementary Haar-cap bound gives this cap mass
at least $\rho^3/24$. Hence
$e^{-r(t)}\ge e^{-1/2}/[24(1+t)^{3/2}]$, proving \eqref{eq:rbound}.
The logarithmic derivative $F'/F$ is the expectation of $\cos\theta$
under its positive tilted angular law. It is strictly less than one.
Pairing angles $\theta$ and $\pi-\theta$ shows $F'(t)>0$ for $t>0$.
\end{proof}
Define
\begin{equation}
 \mathcal E_D(z)=E_0(z)+\sum_{e\in D}m_e(z)
               =\min_{U_D}E(U_D,z),\qquad
 W_D(z)=-\sum_{e\in D}r(Hb_e(z)).
 \label{eq:effective}
\end{equation}
Then the exact same density is
\begin{equation}
 \nu_D(\dd z)=Z_H^{-1}e^{-H\mathcal E_D(z)+W_D(z)}\dd z.
 \label{eq:split}
\end{equation}
This separates an order-$H$ leading action from a logarithmic remainder;
it does not drop the boundary-forced minimum as a scalar.

\paragraph{Why the separation matters.}
For one actual frozen star, both $b=0$ (three $+I$ and three $-I$
staples) and $b=6$ (identical staples) are realizable. The log-integral
factor $-6H+\log F(Hb)$ thus has oscillation
\[
 \log F(6H)=6H-r(6H)\ge6H-r_H.
\]
It is not an $O(\log H)$ perturbation of $E_0$ uniformly in all boundary
data. Only the remainder after extracting $6-b$ has that bound.
Moreover $r(H|B|)$ and $6-|B|$ separately have a cusp at $B=0$;
the full expression does not. Indeed
$F(t)=1+t^2/8+O(t^4)$, so
$-H(6-b)-r(Hb)=-6H+\log F(Hb)$ is smooth in the four coordinates of
$B$. A Hessian argument must not differentiate the separated leading
minimum as if it were everywhere smooth.

\section{Support spreading and inherited reflection positivity}
\begin{proposition}[Failure of elementary-cell closure]\label{prop:nonclosure}
For $H>0$, even the marginal from one removed link generally cannot be
written, up to a constant, as the exponential of a sum of functions
supported on individual original elementary four-cells of retained links.
This statement concerns the original cells and no new auxiliary variables.
\end{proposition}
\begin{proof}
Take $e=(0,e_0)$. Set retained links to identity except the two parallel
middle links of its $+e_1$ and $-e_1$ staples, with values
$\exp(i\alpha\sigma_3)$ and $\exp(i\beta\sigma_3)$ respectively. These
two links lie at transverse coordinates $+1$ and $-1$; no original
unit cell contains both, including for period four. The other four
staples are identity. In quaternion coordinates,
\[
 b(\alpha,\beta)=|4+e^{i\alpha}+e^{i\beta}|,\quad
 b^2=18+8\cos\alpha+8\cos\beta+2\cos(\alpha-\beta).
\]
At the origin $b=6$, its first derivatives vanish, and
$\partial_\alpha\partial_\beta b=1/6$. Consequently
\begin{equation}
 \left.\partial_\alpha\partial_\beta\log F(Hb)\right|_{0,0}
 =\frac H6\frac{F'(6H)}{F(6H)}>0.
 \label{eq:mixed}
\end{equation}
Every cell-supported summand, and every retained original plaquette,
depends on at most one of these two variables after the others are fixed.
Its mixed four-point difference is zero. The smooth exact log-density
has a nonzero mixed difference for sufficiently small equal nonzero
angles by \eqref{eq:mixed}. This argument does not assume differentiability
of a putative cell decomposition. If equality is asserted only Haar-almost
everywhere, Fubini makes its mixed four-point difference zero almost
everywhere in the common exterior and four link values. The nonzero
difference of the continuous exact log-density persists on open
neighborhoods of the displayed choices, which have positive Haar measure.
Thus an almost-everywhere decomposition is ruled out as well.
\end{proof}
Larger supports do not by themselves destroy reflection positivity of
the \emph{complete exact marginal}. If a site reflection $\vartheta$
preserves $D$ as an unoriented link set, the retained-coordinate projection
commutes with it. A bounded real retained positive-half function lifts to a
fine positive-half function, so
\begin{equation}
 \E_{\nu_D}[f\vartheta f]=\E_{\mu_H}[\widetilde f\vartheta\widetilde f]
 \ge0.
 \label{eq:rp}
\end{equation}
This uses the fine model's site-reflection positivity proved in the
preceding notes. It does not assert positivity of the individual
$m_e$ or $r(Hb_e)$ interactions, or of an arbitrary non-invariant marginal.

A concrete nonempty periodic pattern is
\begin{equation}
 D_0=\{(x,x+e_0):x_1,x_2,x_3\text{ even},\ x_0\text{ arbitrary}\}.
 \label{eq:pattern}
\end{equation}
It has $V/8$ links. A plaquette involving direction zero has two parallel
direction-zero links at opposite transverse parities, hence contains at
most one removed link. A plaquette in the other directions contains none. Normal site reflection changes the starting $x_0$ by an odd shift
but $x_0$ is unrestricted; transverse site reflections preserve even
parity. Thus $D_0$ is invariant under every integer-site coordinate
reflection. The marginal retains those positivities, translations by one
in direction zero, and translations by two in the other directions.
It is not a Wilson model on a rescaled complete hypercubic lattice.

\section{Local closure and moment bounds in the exact marginal}
Fix a nonempty $J\subset R$, with $n=|J|$. Let
\begin{equation}
 A_D(J)=\{e\in D:F_e\cap J\ne\varnothing\},\quad a=|A_D(J)|,\quad
 K=J\cup A_D(J),\quad s=|K|=n+a,
 \label{eq:closure}
\end{equation}
where $F_e$ consists of the 18 retained detour links in the removed star.
Since an original plaquette contains at most one removed link and each
retained link belongs to six plaquettes,
\begin{equation}
 a\le6n,\qquad s\le7n.
 \label{eq:closurebound}
\end{equation}
Keep every term of every affected star, not only its plaquettes which
happen to touch $J$. Define the local leading energy
\begin{equation}
 G_J(z)=\sum_{\substack{q\cap J\ne\varnothing\\q\cap D=\varnothing}}s(z_q)
                +\sum_{e\in A_D(J)}m_e(z).
 \label{eq:G}
\end{equation}
Then $\mathcal E_D-G_J$ depends only on retained links outside $J$.
For the local remainder $W_J=-\sum_{e\in A_D(J)}r(Hb_e)$, $W_D-W_J$
also depends only on that exterior. Every $J$-edit obeys
\begin{equation}
 |W_D(z)-W_D(y)|\le a r_H\le6nr_H.
 \label{eq:editcost}
\end{equation}
The full leading energy may change by order $Hn$ in the exponent; that
change is accounted for by its energy gain, not by \eqref{eq:editcost}.

For an original fine-link set $K$, write $E_K$ for the sum of all
original plaquettes touching $K$. No plaquette touching $K$ contains a
removed link outside $A_D(J)$. Thus, as a function of the remaining
retained data and the removed variables $u_A$,
\begin{equation}
 E_K(u_A,z)=
 \sum_{\substack{q\cap J\ne\varnothing\\q\cap D=\varnothing}}s(z_q)
       +\sum_{e\in A_D(J)} E_e(u_e,z),\qquad
 G_J(z)=\min_{u_A}E_K(u_A,z).
 \label{eq:lift}
\end{equation}
Let $w=z_{R\setminus J}$ and
\begin{equation}
 g_J(w)=\min_{z_J}G_J(z_J,w)
       =\min_{z_J,u_A}E_K(u_A,z_J,w).
 \label{eq:gmin}
\end{equation}
Both minima are continuous by compactness. Removed variables outside
$A_D(J)$ are irrelevant to these functions and cannot secretly supply
frozen-boundary terms.

Recall the established original Wilson moment \cite{boundary}:
\[
 \E_{\mu_H}e^{HE_K/128}\le e^{b_H|K|},\qquad
 b_H=\tfrac12a_H+\frac{\log2}{64},\quad
 a_H=\tfrac32+\log24+\tfrac32\log(8H).
\]
\begin{theorem}[Inherited local energy and boundary-minimum tails]
\label{thm:moment}
For the exact marginal of every admissible fixed $D$,
\begin{equation}
 \E_{\nu_D}e^{HG_J/128}\le e^{b_H(n+a)},\qquad
 \E_{\nu_D}e^{Hg_J/128}\le e^{b_H(n+a)}.
 \label{eq:moment}
\end{equation}
Consequently for $v\ge0$,
\begin{equation}
 \nu_D(G_J\ge v),\ \nu_D(g_J>v)
 \ \le\ \min\{1,e^{-Hv/128+b_H(n+a)}\}.
 \label{eq:tails}
\end{equation}
The second is a tail in the \emph{actual retained exterior marginal},
not at an arbitrarily fixed boundary. The estimates need no new
chessboard inequality on the decimated graph.
\end{theorem}
\begin{proof}
For every fine configuration, $0\le g_J\le G_J\le E_K$ by
\eqref{eq:lift}--\eqref{eq:gmin}. Expectation of a retained-coordinate
function under $\nu_D$ equals its expectation under $\mu_H$ exactly.
Apply the original moment and then Markov's inequality. In particular,
this proof works even if $D$ is not reflection invariant.
\end{proof}

\section{A retained-link sublevel kernel without differentiating a cusp}
\begin{theorem}[Retained near-minimum replacement]\label{thm:kernel}
For every $J$ there is a Borel gauge-covariant retained-link probability
kernel, independent of the original values in $J$, freezing all retained
links outside $J$, with conditional density on the edited coordinates
relative to product probability Haar on $\SU^J$ at most
\begin{equation}
 M_J=(24H^{3/2})^n=e^{\ell_H n},\qquad
 \ell_H=\log24+\tfrac32\log H,
 \label{eq:density}
\end{equation}
and every output satisfies
\begin{equation}
 G_J(y_J,w)\le g_J(w)+12n/H.
 \label{eq:kernelenergy}
\end{equation}
The constant uses $n$, not the larger lift size $n+a$.
\end{theorem}
\begin{proof}
Use probability Haar restricted to the retained sublevel
$\mathcal B_J(w)=\{y_J:G_J(y_J,w)\le g_J(w)+12n/H\}$.
It remains to bound its Haar volume.

For a fixed $w$, take any full minimizer $(v_J,v_A)$ of the smooth
fine energy $E_K$ in \eqref{eq:gmin}. Hold $v_A$ fixed and vary only
$v_e\exp(tX_e)$ for $e\in J$, with $\norm{X_e}_{\rm op}\le H^{-1/2}$.
The first derivative of $E_K$ at zero vanishes. As in the earlier
minimizer-cap proof \cite{rough}, the product rule and original
four-link/six-plaquette incidences give
\[
 |E_K''(t)|\le24\sum_{e\in J}\norm{X_e}_{\rm op}^2.
\]
Terms of $E_K$ not touching $J$ have zero derivative. Taylor's integral
remainder therefore gives
$E_K(v_Je^X,v_A,w)\le g_J(w)+12n/H$.
Since $G_J(v_Je^X,w)\le E_K(v_Je^X,v_A,w)$, the full product of $n$
retained caps lies in $\mathcal B_J(w)$. Its Haar mass is at least
$(24H^{3/2})^{-n}$. This proof differentiates only the smooth lifted
energy, not the potentially nonsmooth $G_J$.

Continuity of $G_J$ and $g_J$ makes the sublevel Borel and its
positive-volume normalization Borel by parameter integration.
Gauge transformations carry the minimized fine configurations and the
retained sublevel into their transformed counterparts, preserving Haar
volume. Normalized restriction is therefore gauge covariant. No
measurable choice of $(v_J,v_A)$ is used to define the kernel.
\end{proof}
This kernel does not resample the removed links into the output, and is
not claimed to preserve a finer conditional distribution, reflection
positivity, mixed-loop holonomies, or a continuum topological charge.
It is an existence result, not an implemented generic minimizer or
sublevel sampler.

\section{Which rough labels survive marginalization}
The local leading energy is a sum over a fixed inventory of nonnegative
\emph{effective labels}:
\begin{itemize}
\item an original plaquette $q$ containing no removed link, with
value $s(z_q)$ and its four retained links as declared support;
\item a removed-link star $e\in D$, with value $m_e(z)$ and its 18
retained detour links $F_e$ as declared support.
\end{itemize}
These are labels of the exact minimized action, not a reassignment of
independent plaquette matrices. Each original removed plaquette has
ceased to be an observable of $z$ alone. In particular, an $m_e$ label
is not the roughness of an individual eliminated plaquette.

Declare a label rough if its value exceeds $p^2/(2H)$, for $p>0$.
Let $N_p^D(J)$ count rough labels whose declared support meets $J$.
The values of all effective labels whose declared support meets $J$ sum
exactly to $G_J$; the sum over only the rough ones is bounded by $G_J$.
An unremoved plaquette is rough
precisely at the original curvature threshold $p/\sqrt H$, but a star
value ranges up to six rather than two. Thus all effective rough labels
are empty at $p^2\ge12H$; the earlier all-plaquette cutoff $p^2\ge4H$
must not be applied to star labels.

For a positive integer $k$ split the event
$\mathcal R_{J,k}=\{N_p^D(J)\ge k\}$ at
\begin{equation}
 \mathcal A_{J,k}=\mathcal R_{J,k}\cap\{g_J\le kp^2/(4H)\},\qquad
 \mathcal B_{J,k}=\{g_J>kp^2/(4H)\}.
 \label{eq:split_events}
\end{equation}
These are events of the exact marginal and its retained boundary.
If $kp^2\ge96n$, sources in $\mathcal A_{J,k}$ and every kernel output
obey
\begin{equation}
 N_p^D(J;y)\le3k/4,\qquad
 H(\mathcal E_D(z)-\mathcal E_D(y))\ge kp^2/8.
 \label{eq:gain}
\end{equation}
Indeed $G_J(z)>kp^2/(2H)$ while
$G_J(y)\le kp^2/(4H)+12n/H$; unaffected leading terms cancel. The
output rough count is at most $2HG_J(y)/p^2$.

\begin{theorem}[One-step weighted repair and exceptional-label control]
\label{thm:repair}
Let $kp^2\ge96n$. For the retained sublevel kernel $\mathcal K_J$,
set $T_{J,k}(z,\dd y)=\one_{\mathcal A_{J,k}}(z)\mathcal K_J(z,\dd y)$.
As measures on the retained link space,
\begin{equation}
 \nu_D T_{J,k}\le
 \exp[-kp^2/8+\ell_Hn+a r_H]\nu_D.
 \label{eq:output}
\end{equation}
The exceptional branch in the \emph{same exact marginal} satisfies
\begin{equation}
 \nu_D(\mathcal B_{J,k})\le
 \min\{1,\exp[-kp^2/512+b_H(n+a)]\}.
 \label{eq:exception}
\end{equation}
Consequently
\begin{equation}
 \nu_D(\mathcal R_{J,k})\le\min\{1,
 e^{-kp^2/8+\ell_Hn+a r_H}
 +e^{-kp^2/512+b_H(n+a)}\}.
 \label{eq:repair}
\end{equation}
No new unknown or restricted partition function enters these bounds.
\end{theorem}
\begin{proof}
On accepted input/output pairs, the exact density \eqref{eq:split},
the gain \eqref{eq:gain}, and the residual edit bound \eqref{eq:editcost}
give
\[
 e^{-H\mathcal E_D(z)+W_D(z)}
 \le e^{-kp^2/8+a r_H}e^{-H\mathcal E_D(y)+W_D(y)}.
\]
Bound the retained-output kernel density by $e^{\ell_Hn}$ and integrate
the original $z_J$ over probability Haar. Tonelli leaves the full
normalized marginal on the output, proving \eqref{eq:output} against
every nonnegative output test function. Its total accepted mass is
$\nu_D(\mathcal A_{J,k})$. Apply \eqref{eq:tails} at $v=kp^2/(4H)$
for \eqref{eq:exception}, and then use
$\mathcal R_{J,k}\subset\mathcal A_{J,k}\cup\mathcal B_{J,k}$.
\end{proof}
For useful explicit absorption it suffices to add
\begin{equation}
 kp^2\ge16(\ell_Hn+a r_H),\qquad
 kp^2\ge1024b_H(n+a).
 \label{eq:absorb}
\end{equation}
The two terms in \eqref{eq:repair} are then at most
$e^{-kp^2/16}$ and $e^{-kp^2/1024}$. Since $a\le6n$ and
$r_H,\ell_H,b_H=O(\log H)$, these are genuinely region-proportional
costs for one exact step, independent of $|D|$ and the ambient volume.
A direct moment estimate is stronger for source probabilities alone:
\begin{equation}
 \nu_D(N_p^D(J)\ge k)
 \le\min\{1,e^{-kp^2/256+b_H(n+a)}\}.
 \label{eq:rough}
\end{equation}
The split theorem's additional information is its accepted \emph{output}
domination. It neither removes all effective rough labels nor restores
the old plaquette inventory.

\section{Geometric counts, overlap and accepted composition}
Each retained link lies in at most six effective supports: every one
is witnessed by one of its six original incident plaquettes, whose
removed link is either absent or unique. Every effective support has
at most 18 links. Hence the adjacency graph in which retained links
share an effective support has maximum degree at most
$6(18-1)=102$.

Let $C_p^D(e_0)$ be the actual component through a fixed retained root of
all retained links incident to a rough effective label, with size zero
if the root is absent. A component of size at least $n$ contains a
connected $n$-link subset $J$ through that root. Each such link has a
rough witness; one label can witness at most 18 links, so
$N_p^D(J)\ge\lceil n/18\rceil$. Counting at most $102^{2(n-1)}$ rooted
connected subsets and applying \eqref{eq:rough} with $n+a\le7n$ gives
\begin{equation}
 \nu_D(|C_p^D(e_0)|\ge n)
 \le\min\{1,102^{2(n-1)}e^{-n[p^2/4608-7b_H]}\}.
 \label{eq:components}
\end{equation}
This pays source-dependent geometry in a single fixed marginal, not by
conditioning on a selected support and then pretending its law is the
original conditional Gibbs law.

Accepted operators satisfying $\nu_DT_i\le c_i\nu_D$ compose by order
of measures: $\nu_DT_1\cdots T_r\le(\prod_i c_i)\nu_D$.
All gates and gains must be evaluated at the actual intermediate state.
For a concrete adaptive accepted kernel, choose the \emph{current}
nonempty component $J=C_p^D(e_0)$, let $k=\lceil|J|/18\rceil$, and
accept $\mathcal K_J$ only on its low-minimum branch. Rejected sources
contribute mass zero to this subprobability kernel. If
\begin{equation}
 p^2\ge1728,\qquad d_H=p^2/144-\ell_H-6r_H>2\log102,
 \label{eq:adaptivebudget}
\end{equation}
then counting the possible $J$ in \eqref{eq:output} gives
\begin{equation}
 \nu_DT\le c_H\nu_D,\qquad
 c_H=\frac{e^{-d_H}}{1-102^2e^{-d_H}},\qquad
 \nu_DT^r\le c_H^r\nu_D.
 \label{eq:adaptive}
\end{equation}
Here $c_H<1$ when $d_H>\log10405$. This concerns repeated \emph{repairs
inside one fixed marginal}, not repeated integrations or an RG flow.
A process that stays put on rejected sources is a different operator.

Retained exceptional identity branches still do not multiply as small
operators: encountering the same event twice costs its probability once,
not its square. For joint labels evaluated on one configuration, one
available estimate is as follows. For a nonempty finite family of nonempty
fixed $J_i$ let
$K_i=J_i\cup A_D(J_i)$, $K_\cup=\bigcup_iK_i$, and
$M=\max_q|\{i:q\cap K_i\ne\varnothing\}|$. Then
\begin{equation}
 \E_{\nu_D}\exp\left[\frac H{128M}\sum_i g_{J_i}\right]
 \le e^{b_H|K_\cup|}.
 \label{eq:joint}
\end{equation}
The proof is the pointwise bound
$\sum_i g_{J_i}\le\sum_i E_{K_i}\le M E_{K_\cup}$ under the fine lift.
Even disjoint retained $J_i$ can have overlapping lifted closures; $M$
must be computed, not set to one. This is not a joint estimate for
labels evaluated in different, adaptively changed exteriors.

\section{What this exact step does not establish}
\paragraph{Not the previous cell class.}
Proposition~\ref{prop:nonclosure} is an actual counterexample to original
cellwise additive support after one removal. Even after the order-$H$
minimum is extracted, the remainder retains larger supports; the leading
energy is a sum of minima, not a Wilson plaquette sum. The proof of the
new bound does not import the fourth-pass cell theorem without checking
its hypotheses: it instead uses the exact fine-to-marginal identity,
the closure $K$, and a smooth fine minimizer.

\paragraph{The factorization hypothesis matters.}
If two removed links share a plaquette, \eqref{eq:decomposition} fails:
the plaquette energy depends on both variables. Multiplying one-link
integrals with the other removed link treated as frozen is not joint
marginalization. The accompanying exact finite-group negative control
exhibits this failure. Disjoint stars can still share retained variables;
their effective weights are not probabilistically independent.

\paragraph{Not an iterated renormalization theorem.}
This operation does not introduce coarse links of a new lattice spacing,
match a running coupling, or control the next exact marginal. The induced
multi-staple factors will generally couple links chosen for a later
removal even if they share no \emph{original} plaquette. The one-step
closure factor seven, residual costs and effective-label inventory cannot
simply be reused at every depth. Starting from the fourth-pass weighted
cell class also need not give a product integral: a cell potential can
couple several otherwise plaquette-disjoint removed links.

\paragraph{No conditional repair of impossible boundaries.}
The bounds on $m_e$ and $g_J$ are marginal costs, not claims that arbitrary
frozen holonomies can be made small. The frustrated star, finite-area
loop constraints and flat winding holonomies from earlier passes remain.
A low leading-energy label does not specify an all-small filling of each
integrated star; a partial reduction in effective rough count is not an
inventory-preserving surgery. Signed or complex covering activities,
restricted normalizers and continuum-sector identifications are not
introduced here.

\paragraph{Next concrete target.}
Specify a second exact elimination or a true block-link map, track the
induced supports and retained defect inventory, and determine whether a
positive endpoint class has depth-uniform local energy moments and edit
costs. A valid obstruction or explicit growth law is as useful as a
closure lemma: do not assume the one-step constant is stable. Joint
retained defects and mandatory signed covering identities remain separate
requirements. The fixed-physical-size defect test, cutoff comparison and
nontrivial continuum reconstruction are still downstream.

\section{Verification scope}
The executable checks test probability-Haar normalization, the Bessel
power series versus independent angular quadrature, the zero-staple
endpoint and logarithmic remainder, exact integer noncommuting quaternion
link algebra, sparse-pattern reflection geometry, lifted supports,
conditional minima in finite-group models, and output/overlap budgets.
Exact $\mathbb Z_2$ and quaternion-group integrations are finite-model
diagnostics, not calculations of the four-dimensional $\SU$ partition
function or formal proofs of the analytic claims. Deliberately corrupted
normalizations, lost star terms and falsely factorized shared plaquettes
are negative controls. Existing checks and Lean scalar certificates are
retained; no new analytic theorem is Lean-formalized. The separate audit
files distinguish independent agent review from the author's checks;
neither is external expert certification or a novelty assessment.

\begin{thebibliography}{9}
\bibitem{dlmf}
NIST, \emph{Digital Library of Mathematical Functions},
modified Bessel functions, equations 10.32.2 and 10.25.2, inspected
October 9, 2026. Exact source locators are in \texttt{sources.md}.
\bibitem{boundary}
\emph{Marginal boundary energy costs in the Wilson model}, third-pass
local note, October 9, 2026,
\texttt{../boundary\_cost/boundary\_cost.tex}, Theorem 1.1.
The input remains the documented written theorem, not a Lean-certified
analytic result.
\bibitem{rough}
\emph{Local link replacement with explicit boundary conditions},
second-pass local note, October 9, 2026,
\texttt{../rough\_region/rough\_region\_replacement.tex}, minimizer-cap
lemma and robust one-link obstruction.
\bibitem{weighted}
\emph{Weighted boundary costs and accepted local repairs}, fourth-pass
local note, October 9, 2026,
\texttt{../weighted\_boundary/weighted\_boundary.tex}, weighted output
domination, elementary-cell class and its stated limits.
\end{thebibliography}
\end{document}
